% F3: Heilmeier answers, version 1. F9: version 2.
%   - Write bullet answers first in appendices/my_material.tex; they must be
%     your own. Then answer each question here in full sentences.
%   - Be specific and quantitative. No jargon in question 1. Cite 2 to 3
%     papers in question 2. Each answer is about the research, not about you.
%   - The cost question is optional.
% F9: before revising, save this file unchanged as sections/heilmeier_v1.tex.
% latexmkrc then writes the diff, and \trackededits{heilmeier_v1} shows it.

\section{The Heilmeier Questions}

\begingroup
\renewcommand{\labelenumi}{\textbf{Q\arabic{enumi}.}}
\begin{enumerate}
  \item \textbf{What are you trying to do? Articulate your objectives using
    absolutely no jargon.}\par
    \textit{[Answer.]}
  \item \textbf{How is it done today, and what are the limits of current
    practice?}\par
    \textit{[Answer, with 2 to 3 citations.]}
  \item \textbf{What is new in your approach, and why do you think it will be
    successful?}\par
    \textit{[Answer.]}
  \item \textbf{Who cares? If you are successful, what difference will it
    make?}\par
    \textit{[Answer.]}
  \item \textbf{What are the risks?}\par
    \textit{[Answer.]}
  \item \textbf{How much will it cost?} (optional)\par
    \textit{[Answer or delete this item.]}
  \item \textbf{How long will it take?}\par
    \textit{[Answer.]}
  \item \textbf{What are the midterm and final ``exams'' to check for
    success?}\par
    \textit{[Answer with measurable criteria.]}
\end{enumerate}
\endgroup
